\par\clearpage\SourceLeaf{277}\ArithmeticHead{192}{L}{L. SCHLESINGER, ÜBER GAUSS' FUNKTIONENTHEORETISCHE ARBEITEN.}
\begin{center}\textls[100]{c) Konforme Abbildung der Ellipse. Pentagramma mirificum.}\end{center}
\noindent\hspace*{\parindent}Das Handbuch 19, Be enthält auf den Seiten 172—177 und 222—226
die Lösung des Problems, eine Ellipse konform auf den Einheitskreis abzubilden.
Die beiden getrennten Aufzeichnungen (abgedruckt Werke X, 1,
S. 311—320, vergl. die Bemerkungen S. 323—325) sind zu verschiedenen
Zeiten entstanden. Die erste stammt wohl aus dem Jahre 1834, da auf S. 184
des Handbuchs, bei der (Werke V, S. 609, Art. 7 abgedruckten) Notiz \textit{Das
Induktionsgesetz} die Angabe steht: »Gefunden 1835, Jan. 23. Morgens 7 Uhr,
vor dem Aufstehen«; sie trägt die Überschrift: \textit{Anziehung in der Ebene, umgekehrt
dem Abstand proportional} und beginnt mit der Aufgabe, für eine
gegebene geschlossene Kurve \(L\) alle diejenigen andern \(L'\) zu finden, so dass
der zwischen \(L\) und \(L'\) gelegene Flächenring »keine Anziehung auf den
inneren Raum bewirke«. Der Lösungsversuch ist durchstrichen, bemerkenswert
ist jedoch, dass er die Lösung darauf zurückführen will, wenn \(x,y\) die
Koordinaten der Punkte von \(L\) bedeuten, \(T=x+yi\) »in Gestalt einer Funktion
\(T=f(t)\)«, wo \(t=e^{\varphi i}\) ist, darzustellen, d. h. also auf die Abbildung der
Kurve \(L\) auf den Einheitskreis. — Diese letztere Aufgabe wird nun für den
Fall in Angriff genommen, wo \(L\) eine Ellipse mit den Halbachsen \(a,b\) ist,
und zwar nach der seit frühester Jugendzeit geübten Methode des formalen
Reihenansatzes, indem
\[
\begin{aligned}
x+yi=T&=t+At^3+Bt^5+\ldots,\\
x-yi=T'&=t^{-1}+At^{-3}+Bt^{-5}+\ldots
\end{aligned}
\]
gesetzt und versucht wird, die reellen Koeffizienten \(A,B,\ldots\) so zu bestimmen,
dass die Gleichung der Ellipse
\[
\tfrac12(a^2-b^2)(T^2+T'^2)=TT'(a^2+b^2)-2a^2b^2
\]
befriedigt wird. Die Aufzeichnung bricht nach vergeblichen Versuchen, das
Koeffizientengesetz zu finden ab, darunter ist später die Bemerkung gesetzt:
»Gesetz gefunden 1839, Okt. 9. siehe unten Seite 222«. Die auf dieser Seite
des Handbuchs beginnende zweite Aufzeichnung, die damit genau datiert ist,
nimmt — wohl unter dem Einfluss der in jene Zeit fallenden Untersuchungen
über die nach dem \textsc{Newtonschen} Gesetze wirkenden Kräfte\footnote{Siehe die Abhandlung \textit{Allgemeine Lehrsätze usw.} vom 9. März 1840, Werke V, S. 194, \textit{Anzeige} S. 305.} — zum Aus-\par
\clearpage\SourceLeaf{278}\ArithmeticHead{193}{R}{IXc. KONFORME ABBILDUNG DER ELLIPSE.}
\noindent gangspunkt die Konstruktion der sogenannten \textsc{Greenschen} Funktion der Ellipse
mit dem Mittelpunkt als Pol, indem nämlich (Art. 5, Werke X, 1, S. 315,
316) die Masse \(2\pi A\) so auf die Peripherie der Ellipse verteilt wird, dass ihr
Potential\footnote{Bemerkenswert ist hier das Auftreten der Bezeichnung \textls[100]{Potential}, eine Bezeichnung, die \textsc{Gauss} für den dreidimensionalen Raum im Art. 3 der \textit{Allgemeinen Lehrsätze}, Werke V, S. 200 einführt. Über das Historische hierzu vergl. \textsc{Schaefer}, Werke XI, 2, Abh. 2, S. 95, 96.} in den Punkten der Ellipse dasselbe ist, wie das derselben Masse,
wenn sie im Mittelpunkt der Ellipse konzentriert gedacht wird. Die Differenz
dieser beiden Potentiale ist die \textsc{Greensche} Funktion. \textsc{Gauss} ergänzt
diese Differenz durch Hinzufügen des »imaginären Teils« zum »vollständigen
Potential«, d. h. zu der monogenen Funktion von \(x+yi\), und erhält auf
diese Weise direkt den Logarithmus der Funktion \(t\), die die konforme Abbildung
der Ellipsenscheibe auf die Scheibe des Einheitskreises bewirkt. Wenn
die Halbachsen der Ellipse gleich \(1+\lambda,\ 1-\lambda\) genommen werden, setzt er
\[
T=x+yi=u+\lambda u^{-1},
\]
und findet (Art. 6, Gl. (17), S. 317) für \(t\) die Darstellung
\[
t=\frac{\vartheta_{00}(v|\lambda^2)}{\vartheta_{01}(v|\lambda^2)},
\]
wo wir \(2\pi iv=\log(\lambda u^2)\) gesetzt haben. Für
\[
k=\frac{r(\lambda^2)^2}{p(\lambda^2)^2},\quad \tau=\frac tk
\]
ergibt sich dann (Art. 8, S. 319)
\[
\frac1{p(\lambda^2)^2}\int_0^\tau\frac{d\tau}{\sqrt{(1-\tau^2)(1-k^2\tau^2)}}
=\int_0^T\frac{dT}{\sqrt{4\lambda-T^2}}=i\log u,
\]
so dass in der Bezeichnung von \textsc{Jacobi} (vergl. die Fussnote \textsuperscript{115)})
\[
t=\sqrt{k}\,\sin\operatorname{am}(p(\lambda^2)^2i\log u)
\]
gefunden wird, was mit der von \textsc{H. A. Schwarz} im Jahre 1870\footnote{Siehe \textsc{H. A. Schwarz}, \textit{Gesammelte Abhandlungen} II, S. 102.} gegebenen
Lösung des in Rede stehenden Problems übereinstimmt.
\par
Die Art, wie \textsc{Gauss} die \textsc{Greensche} Funktion der Ellipse bestimmt, schliesst
sich ganz an die in den \textit{Allgemeinen Lehrsätzen}, Art. 29ff. gegebene Formulierung
des allgemeinen Randwertproblems an, namentlich an die des Ar-\par
\HistoricalFoot{X 2 Abh. 2.}{25}
\clearpage\SourceLeaf{279}\ArithmeticHead{194}{L}{L. SCHLESINGER, ÜBER GAUSS' FUNKTIONENTHEORETISCHE ARBEITEN.}
\noindent tikels 36\footnote{Für den hier in Betracht kommenden Fall lautet diese Formulierung (Werke V, S. 240, 241) so: Anstatt einer beliebigen Massenverteilung \(D\), die auf das Innere einer geschlossenen Fläche \(S\) beschränkt ist, lässt sich eine Massenverteilung \(E\) bloss auf der Fläche selbst substituieren mit dem Erfolge, dass die Wirkung von \(E\) der Wirkung von \(D\) gleich wird in allen Punkten des äusseren Raumes. Dazu ist nur erforderlich, dass das Potential von \(D\) mit dem von \(E\) auf der ganzen Fläche \(S\) übereinstimmt, und es wird dann die ganze Masse von \(E\) der Masse von \(D\) gleich sein.}, von der \textsc{Gauss} ausdrücklich bemerkt, dass der in der \textit{Intensitas}, Art. 2
(Werke V, S. 87) und auch an verschiedenen Stellen der \textit{Allgemeinen Theorie
des Erdmagnetismus} (1838, Werke V, S. 119ff.) angeführte Lehrsatz als spezieller
Fall des hier bewiesenen erscheint. Wenn man also hiernach annehmen
darf (vergl. \textsc{Schaefer}, a. a. O., S. 102), dass diese Formulierung der Existentialsätze
des Randwertproblems \textsc{Gauss} schon seit 1831 bekannt war, so zeigt andererseits
der zwischen den beiden Teilen der eben besprochenen Arbeit liegende
Zeitraum von fünf Jahren, dass \textsc{Gauss} den Zusammenhang, der zwischen dem
Randwertproblem der Ebene und der Theorie der konformen Abbildung besteht,
erst 1839 durchschaut hat.
\par
Auf eine genaue Analyse der \textit{Allgemeinen Lehrsätze} können wir hier
nicht eingehen. Es möge nur noch hervorgehoben werden, dass \textsc{Gauss} in
den Artikeln 30, 31 dieser Arbeit (Werke V, S. 232, 233) den Beweis für
die Lösbarkeit des \textls[100]{Randwertproblems} für das Innere eines beliebig gestalteten
Körpers in der Weise zu geben versucht, dass er, ohne dabei seines
vierten Einwandes gegen \textsc{d’Alembert} (Werke III, S. 10, vergl. oben S. 58,
59) zu gedenken, die Existenz einer Massenverteilung auf der Oberfläche als
evident ansieht, für die das über diese Oberfläche erstreckte Integral
\[
\int(V-2U)m\,ds
\]
ein Minimum wird. Darin bedeutet \(V\) das Potential der Massenverteilung
mit der positiven Oberflächendichtigkeit \(m\) und \(U\) eine gegebene endliche und
stetige Funktion auf der Oberfläche. Wir erkennen hier die Schlussweise,
die \textsc{Riemann} (1851) zur Grundlage für die Beweise seiner allgemeinen Existenztheoreme
gemacht und (1857), unter Bezugnahme auf \textsc{Gauss}, als \textls[100]{\textsc{Dirichletsches}
Prinzip} bezeichnet hat, dieselbe, die auch \textsc{Lord Kelvin} (1847) bei
der Lösung der sogenannten zweiten Randwertaufgabe angewandt hat. Die
wechselvollen Schicksale dieser Schlussweise, von ihrer Kritik durch \textsc{Weierstrass}
bis zu ihrer Wiederbelebung durch \textsc{Hilbert}, sind bekannt.
\par\clearpage\SourceLeaf{280}\ArithmeticHead{195}{R}{IXc. PENTAGRAMMA MIRIFICUM.}
\noindent\hspace*{\parindent}Nur wenige Seiten hinter dem ersten Teile der Aufzeichnung über die
Abbildung der Ellipse und aus derselben Zeit (1834) wie diese stammend,
beginnt auf S. 181 des Handbuches 19, Be die Behandlung eines andern geometrischen
Problems, dessen Lösung gleichfalls mit der Theorie der elliptischen
Funktionen zusammenhängt, nämlich die des von \textsc{Gauss} als Pentagramma
mirificum bezeichneten sphärischen Fünfecks, dessen Diagonalen Quadranten
sind. Die Untersuchung wird dann auf den Seiten 240—251 des Handbuchs
fortgesetzt, auf S. 249 steht das Datum »1843 April 20.«; einige weitere
hierher gehörige Aufzeichnungen finden sich auf einzelnen Zetteln. Abgedruckt
sind diese Stücke Werke III, S. 481—490 und VIII, S. 106—111. Die erste
Aufzeichnung von 1834 (Werke III, S. 481—484, Zeichnung und Art. 1, 2)
erwähnt den Zusammenhang mit den elliptischen Funktionen ebensowenig, wie
die erste Notiz über die Abbildung der Ellipse. Vielmehr wird es durch die
Zeichnung auf S. 481 und auch durch die Formeln des Art. 1, S. 482, die
sich direkt auf die bei der Figur des Pentagramma auftretenden rechtwinkligen
sphärischen Dreiecke beziehen, wahrscheinlich gemacht, dass \textsc{Gauss} an
die Untersuchungen in \textsc{J. Napiers} \textit{Mirifici Logarithmorum canonis descriptio} (Edinburgh
1614), lib. II, cap. IV\footnote{Ausführlich besprochen von \textsc{Fricke}, Werke VIII, S. 112, 113. \textsc{Gauss} dürfte auch \textsc{J. H. Lamberts} \textit{Beiträge zum Gebrauche der Mathematik} I, 1765, 2. Auflage 1792, gekannt haben, wo in der dritten Abhandlung: \textit{Anmerkungen und Zusätze zur Trigonometrie}, § 13—31 das Pentagramma behandelt wird. Weitere Literatur siehe bei \textsc{Zacharias}, Encyklopädie, Band III, 1, II, S. 1045.} angeknüpft hat; die Bezeichnung als Pentagramma
\textls[100]{mirificum} wäre dann auch in Anlehnung an den Titel des genannten Werkes
gewählt. Ähnlich wie für die Abbildung der Ellipse, erkennt \textsc{Gauss} auch
hier die Beziehung zu den elliptischen Funktionen erst mehrere Jahre später,
nämlich in der Aufzeichnung S. 247—249 des Handbuches (Werke III, S. 488
—490, Art. 6—8), die mit den (im Abdruck weggelassenen) Worten beginnt:
»Quintessenz aller Untersuchungen, mit zum Teil abgeänderten Bezeichnungen«\footnote{\textsc{Fricke} vermutet, dass \textsc{Gauss} die Beziehung zu den elliptischen Funktionen »auf Anregung von \textsc{Jacobis} Abhandlung \textit{Über die Anwendung der elliptischen Transzendenten auf ein bekanntes Problem der Elementargeometrie} (1828, \textsc{Jacobis} Werke I, S. 277) erkannt hat«. Ein direkter Hinweis auf den Zusammenhang zwischen Pentagramma und dem Schliessungsproblem findet sich bei \textsc{Gauss} nicht; sein Hinweis im Art. 8, Werke III, S. 490: »in der Bedeutung von \textsc{Jacobi} p. 31« bezieht sich auch nicht auf die genannte Abhandlung \textsc{Jacobis} (die \textsc{Crelles} Journal 3, S. 376ff. steht), sondern auf dessen \textit{Fundamenta nova} (1829), wo in der Tat auf Seite 31 (S. 82 des I. Bandes von \textsc{Jacobis} Werken) die von \textsc{Gauss} a. a. O. benutzten Zeichen \(K,k,\Delta\varphi\) eingeführt werden. Die Vermutung \textsc{Frickes} scheint hiernach wenig glaubhaft, zumal auch kein Anzeichen dafür spricht, dass \textsc{Gauss} jene Abhandlung \textsc{Jacobis} näher angesehen hat. Wir bemerken noch, dass \textsc{Gauss} das Pentagramma 1842 als Preisaufgabe der philosophischen Fakultät gestellt hat, siehe Werke XII, S. 221.}.\par
\HistoricalFoot{}{25\textsuperscript{*}}
\clearpage\SourceLeaf{281}\ArithmeticHead{196}{L}{L. SCHLESINGER, ÜBER GAUSS' FUNKTIONENTHEORETISCHE ARBEITEN.}
\noindent Diese Beziehung kann in folgender Weise dargestellt werden\footnote{Vergl. \textsc{L. Schlesinger}, \textsc{Crelles} Journal 124, 1901, S. 36ff.; die geometrischen Artikel der Gaussschen Aufzeichnungen Werke III, Art. 1—5 und VIII, Art. 9—11 sind von \textsc{Fricke}, Werke VIII, S. 112—117 ausführlich erläutert worden.}. Wir nehmen
auf der Kugel vom Halbmesser 1 ein Koordinatensystem \(p,q\), für das das
Linienelement die Form \(2(p^2+q^2+1)^{-1}(dp^2+dq^2)^{\frac12}\) erhält, und deuten
\[
x=2p(1-p^2-q^2)^{-1},\quad y=2q(1-p^2-q^2)^{-1}
\]
als rechtwinklige Koordinaten in einer Ebene \(E\). Sind dann \(p_k,q_k\) \((k=1,2,\ldots,5)\)
die Koordinaten der Eckpunkte des Pentagramms, \(x_k,y_k\) die entsprechenden
Werte von \(x,y\), so stellt sich die Bedingung, dass die fünf Diagonalen Quadranten
sind, in der Form dar:
\begingroup\makeatletter\tagsleft@true\let\veqno\leqno\makeatother
\begin{equation*}
x_kx_{k+2}+y_ky_{k+2}+1=0,\tag*{(1)}
\end{equation*}
\endgroup
wo die Indizes modulo 5 zu reduzieren sind\footnote{Vergl. die Gln. (1) bei \textsc{Fricke}, Werke VIII, S. 114.}. Legen wir also in der Ebene
\(E\) durch die 5 Punkte \(x_k,y_k\) einen Kegelschnitt \(C_1\) und konstruieren dessen
Polarfigur \(C_2\) in bezug auf den imaginären Kreis \(x^2+y^2=-1\), so ist das von
diesen fünf Punkten gebildete Fünfeck \(C_1\) eingeschrieben und \(C_2\) umgeschrieben,
womit die Beziehung zwischen dem Pentagramma und einem \textsc{Ponceletschen}
Schliessungsproblem für den Fall des Fünfecks in Evidenz gesetzt ist\footnote{Vergl. \textsc{Gundelfinger}, \textit{Vorlesungen aus der analyt. Geometrie}, 1895, S. 423.}. Die
Kegelschnitte \(C_1,C_2\) werden simultan auf ihre Hauptachsen transformiert; es
seien die so transformierten Gleichungen
\begingroup\makeatletter\tagsleft@true\let\veqno\leqno\makeatother
\begin{equation*}
\begin{gathered}
G'u^2+G''v^2+G=0,\\
GG''u^2+GG'v^2+G'G''=0,
\end{gathered}\tag*{(2)}
\end{equation*}
\endgroup
wo in Übereinstimmung mit \textsc{Gauss} (Art. 5, Werke III, S. 487) \(G<0,\ G'>0,\)
\(G''>0\) sein möge. Führt man durch die Gleichungen
\[
u=\sqrt{-\frac G{G'}}\cos\varphi,\quad v=\sqrt{-\frac G{G''}}\sin\varphi
\]
die exzentrische Anomalie für \(C_1\) ein (Art. 7, a. a. O. S. 489) und bezeichnet
ihre Werte für die 5 Punkte \((x_k,y_k)\) mit \(\varphi_k\), so lauten die den Gleichungen
(1) entsprechenden Relationen in den transformierten Koordinaten
\begingroup\makeatletter\tagsleft@true\let\veqno\leqno\makeatother
\begin{equation*}
\begin{gathered}
GG''\cos\varphi_k\cos\varphi_{k+2}+GG'\sin\varphi_k\sin\varphi_{k+2}-G'G''=0\\
(k=1,2,\ldots,5).
\end{gathered}\tag*{(1\,a)}
\end{equation*}
\endgroup
\par\clearpage\SourceLeaf{282}\ArithmeticHead{197}{R}{IXc. PENTAGRAMMA MIRIFICUM.}
\noindent Durch Verbindung zweier dieser Gleichungen ergeben sich dann die von
\textsc{Gauss} (Art. 7, a. a. O. S. 489) aufgestellten Relationen
\begingroup\makeatletter\tagsleft@true\let\veqno\leqno\makeatother
\begin{equation*}
\begin{gathered}
\frac{\sin\frac12(\varphi_k+\varphi_{k-1})}{\cos\frac12(\varphi_k-\varphi_{k-1})}=\frac G{G''}\sin\varphi_{k+2},\qquad
\frac{\cos\frac12(\varphi_k+\varphi_{k-1})}{\cos\frac12(\varphi_k-\varphi_{k-1})}=\frac G{G'}\cos\varphi_{k+2},
\end{gathered}\tag*{(3)}
\end{equation*}
\endgroup
die sofort an die Multiplikationsformeln der elliptischen Funktionen erinnern.
Setzt man nun mit \textsc{Gauss} (Art. 8, a. a. O. S. 490)
\[
k^2=\frac{\dfrac1{G'^2}-\dfrac1{G''^2}}{\dfrac1{G'^2}-\dfrac1{G^2}},\quad
z=\int_0^\varphi\frac{d\varphi}{\sqrt{1-k^2\sin^2\varphi}}
\]
und mit den \textsc{Jacobischen} Bezeichnungen
\[
\begin{gathered}
\varphi=\operatorname{am}z,\quad\varphi_k=\operatorname{am}z_k,\\
\frac{G'^2}{G^2}=\cos^2\operatorname{am}\tau,\quad
\frac{G'^2}{G''^2}=1-k^2\sin^2\operatorname{am}\tau=\Delta^2\operatorname{am}\tau,
\end{gathered}
\]
so ergibt die Division der Gleichungen (3) durcheinander
\[
\operatorname{tg}\tfrac12(\operatorname{am}z_k+\operatorname{am}z_{k-1})
=\Delta\operatorname{am}\tau\mathbin{.}\operatorname{tg}\operatorname{am}z_{k+2}.
\]
Vergleicht man diese Gleichung, wie es \textsc{Jacobi} beim Schliessungsproblem tut\footnote{Siehe die in der Fussnote \textsuperscript{304)} genannte Abhandlung, \textsc{Jacobis} Werke I, S. 285.},
mit der aus der Theorie der elliptischen Funktionen bekannten Formel
\[
\operatorname{tg}\tfrac12(\operatorname{am}z+\operatorname{am}(z+2\tau))
=\Delta\operatorname{am}\tau\mathbin{.}\operatorname{tg}\operatorname{am}(z+\tau),
\]
so erhält man
\[
\varphi_{k-1}=\operatorname{am}(z_k+2\tau),\quad\varphi_{k+2}=\operatorname{am}(z_k+\tau),
\]
also abgesehen von Vielfachen von \(2\pi\)
\[
\varphi_{2k+1}=\operatorname{am}(z_1+k\tau).
\]
Ist das Fünfeck ein konvexes, so dass beim Durchlaufen der Ecken in der
Reihenfolge 1, 3, 5, 2, 4, 1 die exzentrische Anomalie um \(4\pi\) zunimmt, so ist
\[
\operatorname{am}(z_1+5\tau)=\operatorname{am}z_1+4\pi,
\]
also \(\tau=\frac85K\), so dass sich in Übereinstimmung mit den Formeln von \textsc{Gauss}
(Art. 8, a. a. O. S. 490)
\[
\varphi_k=\operatorname{am}\bigl(z_k+(k-1)\tfrac45K\bigr)\qquad(k=1,2,\ldots,5)
\]
ergibt.
\par\clearpage\SourceLeaf{283}\ArithmeticHead{198}{L}{L. SCHLESINGER, ÜBER GAUSS' FUNKTIONENTHEORETISCHE ARBEITEN.}
\noindent\hspace*{\parindent}Diese Aufzeichnungen von 1843 sind anscheinend die letzten, die \textsc{Gauss}
über elliptische Funktionen gemacht hat. Da er 1791 begonnen hat, sich
mit der Theorie des agM. zu beschäftigen, so erstreckt sich seine Arbeit an
der Theorie dieser Funktionen auf einen Zeitraum von mehr als fünfzig Jahren.
\par
Entspringend aus den beiden Wurzeln des agM. (1791) und der Theorie
der lemniskatischen Funktionen (1797), erwächst der Stamm der Theorie der\linebreak
\textls[100]{elliptischen Funktionen} in den grossen Entdeckungen von November
1799 bis Mitte des Jahres 1800 zu seiner vollen Höhe. Mit Hilfe der durch
das agM. definierten Perioden wird der Sinus lemniscaticus zum Sinus lemniscaticus
in allgemeiner Fassung ausgebaut, und mit der Einführung der vier
Thetafunktionen, die als unendliche Reihen und Produkte dargestellt werden,
gewinnt \textsc{Gauss} das Instrument, das ihm, im Verein mit seiner Reduktionstheorie
der quadratischen Formen, die Herrschaft über die elliptischen Funktionen
und die Modulfunktion sichert.
\par
An die Stelle des agM. tritt dann, wohl um 1801, jedenfalls vor 1806,
die allgemeinere Betrachtung der durch die Reihe \(F(\alpha,\beta,\gamma,x)\) definierten Funktionen,
mit der \textsc{Gauss} in die Theorie der \textls[100]{homogenen linearen Differentialgleichungen
zweiter Ordnung} einlenkt und die Grundlagen für
die funktionentheoretische Behandlung derselben schafft.
\par
In welchem Masse er auch als Vorläufer für die Theorie der \textls[100]{automorphen
Funktionen} angesehen werden muss, zeigt nicht nur seine Konstruktion
des Fundamentalbereichs für die Modulfunktion, sondern auch seine
Deutung der linearen Substitution
\[
t'=\frac{(\alpha+\beta i)t-(\gamma+\delta i)}{(\gamma-\delta i)t+\alpha-\beta i}
\]
als Drehung der Kugelfläche in sich selbst, die in zwei, Werke VIII, S. 354
—356 abgedruckten Aufzeichnungen erörtert wird, und das auf einen Fall
transzendenter Dreiecksfunktionen, der sich nicht der Modulfunktion unterordnet,
bezügliche Polygonnetz, das Werke VIII, S. 104, nach einer Zeichnung
auf einem besonderen Blatte, wiedergegeben ist.
